% Options for packages loaded elsewhere
\PassOptionsToPackage{unicode}{hyperref}
\PassOptionsToPackage{hyphens}{url}
\documentclass[
  11pt,
  a4paper,
]{article}
\usepackage{xcolor}
\usepackage[margin=18mm]{geometry}
\usepackage{amsmath,amssymb}
\setcounter{secnumdepth}{-\maxdimen} % remove section numbering
\usepackage{iftex}
\ifPDFTeX
  \usepackage[T1]{fontenc}
  \usepackage[utf8]{inputenc}
  \usepackage{textcomp} % provide euro and other symbols
\else % if luatex or xetex
  \usepackage{unicode-math} % this also loads fontspec
  \defaultfontfeatures{Scale=MatchLowercase}
  \defaultfontfeatures[\rmfamily]{Ligatures=TeX,Scale=1}
\fi
\ifPDFTeX\else
  % xetex/luatex font selection
    \setmainfont[]{Noto Serif}
    \setsansfont[]{Noto Sans}
    \setmonofont[]{Noto Sans Mono}
  \setmathfont[]{Noto Sans Math}
\fi
% Use upquote if available, for straight quotes in verbatim environments
\IfFileExists{upquote.sty}{\usepackage{upquote}}{}
\IfFileExists{microtype.sty}{% use microtype if available
  \usepackage[]{microtype}
  \UseMicrotypeSet[protrusion]{basicmath} % disable protrusion for tt fonts
}{}
\makeatletter
\@ifundefined{KOMAClassName}{% if non-KOMA class
  \IfFileExists{parskip.sty}{%
    \usepackage{parskip}
  }{% else
    \setlength{\parindent}{0pt}
    \setlength{\parskip}{6pt plus 2pt minus 1pt}}
}{% if KOMA class
  \KOMAoptions{parskip=half}}
\makeatother
\ifLuaTeX
\usepackage[bidi=basic]{babel}
\else
\usepackage[bidi=default]{babel}
\fi
\babelprovide[main,import]{russian}
\ifPDFTeX
\else
\babelfont{rm}[]{Noto Serif}
\fi
% get rid of language-specific shorthands (see #6817):
\let\LanguageShortHands\languageshorthands
\def\languageshorthands#1{}
\setlength{\emergencystretch}{3em} % prevent overfull lines
\providecommand{\tightlist}{%
  \setlength{\itemsep}{0pt}\setlength{\parskip}{0pt}}
\usepackage{bookmark}
\IfFileExists{xurl.sty}{\usepackage{xurl}}{} % add URL line breaks if available
\urlstyle{same}
\hypersetup{
  pdftitle={Билеты по дискретным случайным процессам},
  pdflang={ru-RU},
  hidelinks,
  pdfcreator={LaTeX via pandoc}}

\title{Билеты по дискретным случайным процессам}
\author{}
\date{}

\begin{document}
\maketitle

\section{Билет 22. Марковские переходные функции и
матрицы}\label{ux431ux438ux43bux435ux442-22.-ux43cux430ux440ux43aux43eux432ux441ux43aux438ux435-ux43fux435ux440ux435ux445ux43eux434ux43dux44bux435-ux444ux443ux43dux43aux446ux438ux438-ux438-ux43cux430ux442ux440ux438ux446ux44b}

Источники: Ширяев 2, гл. VIII, §1, §3; Сердобольская, лекция 5.

\subsection{1. Короткий
ответ}\label{ux43aux43eux440ux43eux442ux43aux438ux439-ux43eux442ux432ux435ux442}

Переходная функция, или марковское ядро, задает условный закон
следующего состояния при известном текущем состоянии. Если пространство
состояний \((E,\mathcal E)\) произвольное, то переходная функция имеет
вид

\[
P(x,A)=P\{X_{n+1}\in A\mid X_n=x\},
\qquad
x\in E,\ A\in\mathcal E.
\]

В общем пространстве состояний это понимают как регулярную версию
условного распределения \(P(X_{n+1}\in A\mid X_n)\). В счетном случае
запись через условие \(X_n=x\) безопасна буквально. По \(A\) это
вероятностная мера, по \(x\) - измеримая функция.

Если пространство состояний счетное или конечное, то переходная функция
задается матрицей

\[
P=(p_{ij}),
\qquad
p_{ij}=P\{X_{n+1}=j\mid X_n=i\}.
\]

Матрица переходов стохастическая:

\[
p_{ij}\ge0,
\qquad
\sum_j p_{ij}=1.
\]

При соглашении, что распределения записываются строками, эволюция
распределений задается формулой

\[
\pi_{n+1}=\pi_nP.
\]

Переходы за несколько шагов задаются степенями матрицы:

\[
P^{(n)}=P^n.
\]

Главное свойство - уравнение Чепмена-Колмогорова:

\[
P^{(m+n)}(x,A)
=
\int_E P^{(n)}(x,dy)P^{(m)}(y,A).
\]

В матричном случае:

\[
p_{ij}^{(m+n)}
=
\sum_k p_{ik}^{(n)}p_{kj}^{(m)}.
\]

\subsection{2. Полный
ответ}\label{ux43fux43eux43bux43dux44bux439-ux43eux442ux432ux435ux442}

\subsubsection{22.1. Предмет билета и
контекст}\label{ux43fux440ux435ux434ux43cux435ux442-ux431ux438ux43bux435ux442ux430-ux438-ux43aux43eux43dux442ux435ux43aux441ux442}

Этот билет является технической основой для Билета 23 о марковских
цепях. В билете 23 главный объект - сам процесс с марковским свойством.
В билете 22 главный объект - переходный механизм: как задается условное
распределение будущего состояния через настоящее.

\subsubsection{22.2. Определение переходной
функции}\label{ux43eux43fux440ux435ux434ux435ux43bux435ux43dux438ux435-ux43fux435ux440ux435ux445ux43eux434ux43dux43eux439-ux444ux443ux43dux43aux446ux438ux438}

Пусть

\[
X=(X_0,X_1,X_2,\ldots)
\]

\begin{itemize}
\tightlist
\item
  случайная последовательность со значениями в измеримом пространстве
  состояний
\end{itemize}

\[
(E,\mathcal E).
\]

Марковская идея состоит в том, что для предсказания следующего состояния
достаточно знать текущее. Если процесс однороден по времени, то условный
закон перехода из состояния \(x\) во множество \(A\) не зависит от
номера шага \(n\):

\[
P\{X_{n+1}\in A\mid X_n=x\}=P(x,A).
\]

\subsubsection{22.3. Свойства переходного
ядра}\label{ux441ux432ux43eux439ux441ux442ux432ux430-ux43fux435ux440ux435ux445ux43eux434ux43dux43eux433ux43e-ux44fux434ux440ux430}

Функция

\[
P:E\times\mathcal E\to[0,1]
\]

называется переходной функцией, переходной вероятностью или марковским
ядром. Она должна удовлетворять двум условиям:

\begin{enumerate}
\def\labelenumi{\arabic{enumi}.}
\tightlist
\item
  При фиксированном \(x\in E\) функция
\end{enumerate}

\[
A\mapsto P(x,A)
\]

является вероятностной мерой на \((E,\mathcal E)\).

\begin{enumerate}
\def\labelenumi{\arabic{enumi}.}
\setcounter{enumi}{1}
\tightlist
\item
  При фиксированном \(A\in\mathcal E\) функция
\end{enumerate}

\[
x\mapsto P(x,A)
\]

является \(\mathcal E\)-измеримой.

Первое условие говорит, что из каждого состояния задано полноценное
распределение следующего состояния. Второе условие нужно для
корректности интегралов вида

\[
\int_E P(x,A)\,\mu(dx).
\]

\subsubsection{22.4. Преобразование
распределений}\label{ux43fux440ux435ux43eux431ux440ux430ux437ux43eux432ux430ux43dux438ux435-ux440ux430ux441ux43fux440ux435ux434ux435ux43bux435ux43dux438ux439}

Если \(\mu\) - распределение \(X_n\), то распределение \(X_{n+1}\)
находится так:

\[
\mu P(A)=\int_E P(x,A)\,\mu(dx).
\]

Здесь \(\mu P\) - новая вероятностная мера. Это просто формула полной
вероятности: чтобы попасть в \(A\) на следующем шаге, нужно
просуммировать или проинтегрировать по всем возможным текущим состояниям
\(x\).

\subsubsection{22.5. Матрица переходных вероятностей в счетном
случае}\label{ux43cux430ux442ux440ux438ux446ux430-ux43fux435ux440ux435ux445ux43eux434ux43dux44bux445-ux432ux435ux440ux43eux44fux442ux43dux43eux441ux442ux435ux439-ux432-ux441ux447ux435ux442ux43dux43eux43c-ux441ux43bux443ux447ux430ux435}

В счетном пространстве состояний

\[
E=\{1,2,\ldots\}
\]

переходная функция полностью задается числами

\[
p_{ij}=P(i,\{j\}).
\]

Тогда

\[
P(i,A)=\sum_{j\in A}p_{ij}.
\]

Матрица

\[
P=(p_{ij})_{i,j\in E}
\]

называется матрицей переходных вероятностей. Она является стохастической
по строкам:

\[
p_{ij}\ge0,
\qquad
\sum_{j\in E}p_{ij}=1.
\]

\subsubsection{22.6. Соглашение о векторно-матричном
умножении}\label{ux441ux43eux433ux43bux430ux448ux435ux43dux438ux435-ux43e-ux432ux435ux43aux442ux43eux440ux43dux43e-ux43cux430ux442ux440ux438ux447ux43dux43eux43c-ux443ux43cux43dux43eux436ux435ux43dux438ux438}

Важно: в этом конспекте используется строковое соглашение. Распределение
записывается как строка

\[
\pi_n=(\pi_n(i))_{i\in E},
\]

и тогда

\[
\pi_{n+1}=\pi_nP.
\]

Покомпонентно:

\[
\pi_{n+1}(j)
=
\sum_i \pi_n(i)p_{ij}.
\]

Если кто-то записывает распределения столбцами, то формула меняется на

\[
\pi_{n+1}=P^T\pi_n.
\]

Поэтому нужно всегда фиксировать соглашение. Типичная ошибка - объявить
строки стохастическими, а потом умножать так, как будто стохастические
столбцы.

\subsubsection{22.7. Вероятности перехода за несколько
шагов}\label{ux432ux435ux440ux43eux44fux442ux43dux43eux441ux442ux438-ux43fux435ux440ux435ux445ux43eux434ux430-ux437ux430-ux43dux435ux441ux43aux43eux43bux44cux43aux43e-ux448ux430ux433ux43eux432}

Теперь рассмотрим переходы за несколько шагов. Для однородной цепи
одношаговое ядро \(P\) не зависит от времени. Двухшаговое ядро
определяется формулой

\[
P^{(2)}(x,A)
=
\int_E P(x,dy)P(y,A).
\]

Смысл: сначала система из \(x\) попадает в промежуточное состояние
\(y\), потом из \(y\) попадает в \(A\).

Рекурсивно:

\[
P^{(1)}(x,A)=P(x,A),
\]

а для \(n\ge1\)

\[
P^{(n+1)}(x,A)
=
\int_E P^{(n)}(x,dy)P(y,A).
\]

Для матриц это ровно матричное умножение:

\[
p_{ij}^{(2)}
=
\sum_k p_{ik}p_{kj},
\]

а вообще

\[
P^{(n)}=P^n.
\]

\subsubsection{22.8. Уравнение Чепмена ---
Колмогорова}\label{ux443ux440ux430ux432ux43dux435ux43dux438ux435-ux447ux435ux43fux43cux435ux43dux430-ux43aux43eux43bux43cux43eux433ux43eux440ux43eux432ux430}

Главная формула билета - уравнение Чепмена-Колмогорова. Для однородной
цепи:

\[
P^{(m+n)}(x,A)
=
\int_E P^{(n)}(x,dy)P^{(m)}(y,A).
\]

Оно говорит, что переход за \(m+n\) шагов можно разложить через
промежуточное состояние после \(n\) шагов. В счетном случае интеграл
заменяется суммой:

\[
p_{ij}^{(m+n)}
=
\sum_{k\in E}p_{ik}^{(n)}p_{kj}^{(m)}.
\]

В матричной записи:

\[
P^{(m+n)}=P^{(n)}P^{(m)}.
\]

Так как для однородной цепи \(P^{(n)}=P^n\), это то же самое, что
обычное тождество

\[
P^{n+m}=P^nP^m.
\]

\subsubsection{22.9. Неоднородные по времени
цепи}\label{ux43dux435ux43eux434ux43dux43eux440ux43eux434ux43dux44bux435-ux43fux43e-ux432ux440ux435ux43cux435ux43dux438-ux446ux435ux43fux438}

В неоднородном случае переходный механизм может зависеть от времени.
Тогда вместо одного ядра берут семейство

\[
P_n(x,A)=P\{X_{n+1}\in A\mid X_n=x\}.
\]

Для счетного пространства это семейство матриц

\[
P_n=(p_{ij}(n)).
\]

Тогда распределения эволюционируют так:

\[
\pi_{n+1}=\pi_nP_n,
\]

а переход от времени \(r\) к времени \(s\) задается произведением
соответствующих ядер:

\[
P_{r,s}=P_rP_{r+1}\cdots P_{s-1},
\qquad
r<s.
\]

В однородном случае все \(P_n=P\), поэтому

\[
P_{r,s}=P^{s-r}.
\]

\subsubsection{22.10. Обобщение на непрерывное
время}\label{ux43eux431ux43eux431ux449ux435ux43dux438ux435-ux43dux430-ux43dux435ux43fux440ux435ux440ux44bux432ux43dux43eux435-ux432ux440ux435ux43cux44f}

Можно также встретить непрерывное время. Тогда переходные вероятности
записывают как

\[
p_{ij}(t)=P\{X_{s+t}=j\mid X_s=i\}
\]

в однородном случае, а матрицы перехода образуют семейство

\[
P(t)=(p_{ij}(t)),
\qquad
t\ge0.
\]

Тогда

\[
P(0)=I,
\]

строки \(P(t)\) суммируются в единицу, и уравнение Чепмена-Колмогорова
принимает вид

\[
P(t+s)=P(t)P(s).
\]

Для текущего курса дискретных случайных процессов обычно достаточно
дискретного времени, но полезно понимать, что это та же идея.

\subsubsection{22.11. Начальное распределение и конечномерные
распределения}\label{ux43dux430ux447ux430ux43bux44cux43dux43eux435-ux440ux430ux441ux43fux440ux435ux434ux435ux43bux435ux43dux438ux435-ux438-ux43aux43eux43dux435ux447ux43dux43eux43cux435ux440ux43dux44bux435-ux440ux430ux441ux43fux440ux435ux434ux435ux43bux435ux43dux438ux44f}

Переходная функция не является начальным распределением. Начальное
распределение отвечает на вопрос, где процесс находится в момент \(0\):

\[
\pi_0(i)=P\{X_0=i\}.
\]

Переходная матрица отвечает на вопрос, как процесс переходит из
состояния в состояние:

\[
p_{ij}=P\{X_{n+1}=j\mid X_n=i\}.
\]

Только вместе начальное распределение и переходные вероятности задают
конечномерные распределения цепи. В счетном однородном случае:

\[
P\{X_0=i_0,\ldots,X_n=i_n\}
=
\pi_0(i_0)p_{i_0i_1}p_{i_1i_2}\cdots p_{i_{n-1}i_n}.
\]

Именно эта формула связывает билет с Билетом 15: такие конечномерные
распределения согласованы, потому что строки \(P\) суммируются в
единицу. Значит по теореме Колмогорова можно построить процесс с
заданными \(\pi_0\) и \(P\).

\subsection{3. Основные
определения}\label{ux43eux441ux43dux43eux432ux43dux44bux435-ux43eux43fux440ux435ux434ux435ux43bux435ux43dux438ux44f}

\subsubsection{Пространство
состояний}\label{ux43fux440ux43eux441ux442ux440ux430ux43dux441ux442ux432ux43e-ux441ux43eux441ux442ux43eux44fux43dux438ux439}

Измеримое пространство

\[
(E,\mathcal E),
\]

значения в котором принимает процесс \(X_n\). В конечном или счетном
случае обычно берут

\[
E=\{1,2,\ldots,r\}
\]

или

\[
E=\mathbb Z,\ \mathbb N.
\]

\subsubsection{Переходная
функция}\label{ux43fux435ux440ux435ux445ux43eux434ux43dux430ux44f-ux444ux443ux43dux43aux446ux438ux44f}

Функция

\[
P(x,A),
\qquad
x\in E,\ A\in\mathcal E,
\]

которая задает вероятность перехода из состояния \(x\) в множество \(A\)
за один шаг.

\subsubsection{Марковское
ядро}\label{ux43cux430ux440ux43aux43eux432ux441ux43aux43eux435-ux44fux434ux440ux43e}

То же самое, что переходная функция, но с акцентом на измеримую
структуру:

\begin{itemize}
\tightlist
\item
  \(P(x,\cdot)\) - вероятностная мера;
\item
  \(P(\cdot,A)\) - измеримая функция.
\end{itemize}

\subsubsection{Переходная
вероятность}\label{ux43fux435ux440ux435ux445ux43eux434ux43dux430ux44f-ux432ux435ux440ux43eux44fux442ux43dux43eux441ux442ux44c}

В счетном случае:

\[
p_{ij}=P\{X_{n+1}=j\mid X_n=i\}.
\]

Это вероятность перехода из состояния \(i\) в состояние \(j\) за один
шаг.

\subsubsection{Стохастическая
матрица}\label{ux441ux442ux43eux445ux430ux441ux442ux438ux447ux435ux441ux43aux430ux44f-ux43cux430ux442ux440ux438ux446ux430}

Матрица

\[
P=(p_{ij})
\]

называется стохастической по строкам, если

\[
p_{ij}\ge0,
\qquad
\sum_jp_{ij}=1
\]

для каждой строки \(i\).

\subsubsection{Однородность по
времени}\label{ux43eux434ux43dux43eux440ux43eux434ux43dux43eux441ux442ux44c-ux43fux43e-ux432ux440ux435ux43cux435ux43dux438}

Цепь или переходный механизм называется однородным, если условный закон
перехода не зависит от времени:

\[
P\{X_{n+1}\in A\mid X_n=x\}=P(x,A)
\]

для всех \(n\).

В счетном случае:

\[
P\{X_{n+1}=j\mid X_n=i\}=p_{ij}.
\]

\subsubsection{Неоднородные
переходы}\label{ux43dux435ux43eux434ux43dux43eux440ux43eux434ux43dux44bux435-ux43fux435ux440ux435ux445ux43eux434ux44b}

Если переходы зависят от времени, пишут

\[
P_n(x,A)=P\{X_{n+1}\in A\mid X_n=x\}.
\]

В счетном случае:

\[
p_{ij}(n)=P\{X_{n+1}=j\mid X_n=i\}.
\]

\subsubsection{\texorpdfstring{\(n\)-шаговая переходная
функция}{n-шаговая переходная функция}}\label{n-ux448ux430ux433ux43eux432ux430ux44f-ux43fux435ux440ux435ux445ux43eux434ux43dux430ux44f-ux444ux443ux43dux43aux446ux438ux44f}

Для однородной цепи:

\[
P^{(n)}(x,A)=P\{X_n\in A\mid X_0=x\}.
\]

В счетном случае:

\[
p_{ij}^{(n)}=P\{X_n=j\mid X_0=i\}.
\]

\subsubsection{Начальное
распределение}\label{ux43dux430ux447ux430ux43bux44cux43dux43eux435-ux440ux430ux441ux43fux440ux435ux434ux435ux43bux435ux43dux438ux435}

Распределение состояния в момент \(0\):

\[
\pi_0(A)=P\{X_0\in A\}.
\]

В счетном случае:

\[
\pi_0(i)=P\{X_0=i\}.
\]

\subsection{4. Основные теоремы и
свойства}\label{ux43eux441ux43dux43eux432ux43dux44bux435-ux442ux435ux43eux440ux435ux43cux44b-ux438-ux441ux432ux43eux439ux441ux442ux432ux430}

\subsubsection{1. Стохастичность переходной
матрицы}\label{ux441ux442ux43eux445ux430ux441ux442ux438ux447ux43dux43eux441ux442ux44c-ux43fux435ux440ux435ux445ux43eux434ux43dux43eux439-ux43cux430ux442ux440ux438ux446ux44b}

Для любой переходной матрицы

\[
p_{ij}\ge0,
\qquad
\sum_jp_{ij}=1.
\]

Доказательство: при фиксированном \(i\) числа \(p_{ij}\) являются
вероятностями несовместных событий \(\{X_{n+1}=j\}\), которые образуют
полную группу исходов по следующему состоянию.

\subsubsection{2. Переходная функция переводит распределения в
распределения}\label{ux43fux435ux440ux435ux445ux43eux434ux43dux430ux44f-ux444ux443ux43dux43aux446ux438ux44f-ux43fux435ux440ux435ux432ux43eux434ux438ux442-ux440ux430ux441ux43fux440ux435ux434ux435ux43bux435ux43dux438ux44f-ux432-ux440ux430ux441ux43fux440ux435ux434ux435ux43bux435ux43dux438ux44f}

Если \(\mu\) - вероятностная мера на \((E,\mathcal E)\), то

\[
\mu P(A)=\int_E P(x,A)\,\mu(dx)
\]

тоже вероятностная мера.

Проверка:

\[
\mu P(A)\ge0,
\]

счетная аддитивность следует из счетной аддитивности \(P(x,\cdot)\) и
теоремы Леви или монотонной сходимости, а нормировка:

\[
\mu P(E)=\int_E P(x,E)\,\mu(dx)=\int_E1\,\mu(dx)=1.
\]

\subsubsection{3. Эволюция распределений в счетном
случае}\label{ux44dux432ux43eux43bux44eux446ux438ux44f-ux440ux430ux441ux43fux440ux435ux434ux435ux43bux435ux43dux438ux439-ux432-ux441ux447ux435ux442ux43dux43eux43c-ux441ux43bux443ux447ux430ux435}

Если

\[
\pi_n(i)=P\{X_n=i\},
\]

то

\[
\pi_{n+1}(j)
=
\sum_i\pi_n(i)p_{ij}.
\]

В матричной записи:

\[
\pi_{n+1}=\pi_nP.
\]

Доказательство - формула полной вероятности:

\[
P\{X_{n+1}=j\}
=
\sum_iP\{X_n=i\}P\{X_{n+1}=j\mid X_n=i\}.
\]

\subsubsection{4. Двухшаговые
переходы}\label{ux434ux432ux443ux445ux448ux430ux433ux43eux432ux44bux435-ux43fux435ux440ux435ux445ux43eux434ux44b}

В общем пространстве:

\[
P^{(2)}(x,A)
=
\int_E P(x,dy)P(y,A).
\]

В счетном пространстве:

\[
p_{ij}^{(2)}
=
\sum_k p_{ik}p_{kj}.
\]

Это обычное умножение матриц:

\[
P^{(2)}=P^2.
\]

\subsubsection{\texorpdfstring{5. \(n\)-шаговые
переходы}{5. n-шаговые переходы}}\label{n-ux448ux430ux433ux43eux432ux44bux435-ux43fux435ux440ux435ux445ux43eux434ux44b}

Для однородной цепи:

\[
P^{(n)}=P^n.
\]

В счетном случае:

\[
p_{ij}^{(n)}
=
(P^n)_{ij}.
\]

При этом

\[
\pi_n=\pi_0P^n.
\]

\subsubsection{6. Уравнение
Чепмена-Колмогорова}\label{ux443ux440ux430ux432ux43dux435ux43dux438ux435-ux447ux435ux43fux43cux435ux43dux430-ux43aux43eux43bux43cux43eux433ux43eux440ux43eux432ux430-1}

Для всех \(m,n\ge0\):

\[
P^{(m+n)}(x,A)
=
\int_E P^{(n)}(x,dy)P^{(m)}(y,A).
\]

В счетном случае:

\[
p_{ij}^{(m+n)}
=
\sum_kp_{ik}^{(n)}p_{kj}^{(m)}.
\]

В матричном виде:

\[
P^{(m+n)}=P^{(n)}P^{(m)}.
\]

Смысл: переход за \(m+n\) шагов раскладывается по промежуточному
состоянию через \(n\) шагов.

\subsubsection{7. Формула конечномерных
распределений}\label{ux444ux43eux440ux43cux443ux43bux430-ux43aux43eux43dux435ux447ux43dux43eux43cux435ux440ux43dux44bux445-ux440ux430ux441ux43fux440ux435ux434ux435ux43bux435ux43dux438ux439}

Для однородной цепи со счетным пространством состояний:

\[
P\{X_0=i_0,\ldots,X_n=i_n\}
=
\pi_0(i_0)p_{i_0i_1}\cdots p_{i_{n-1}i_n}.
\]

Эта формула показывает, что начальное распределение и переходная матрица
полностью задают закон цепи.

\subsubsection{8. Непрерывное время: матрицы
перехода}\label{ux43dux435ux43fux440ux435ux440ux44bux432ux43dux43eux435-ux432ux440ux435ux43cux44f-ux43cux430ux442ux440ux438ux446ux44b-ux43fux435ux440ux435ux445ux43eux434ux430}

Если рассматривается однородный марковский процесс с конечным числом
состояний и непрерывным временем, то

\[
P(t)=(p_{ij}(t)),
\qquad
p_{ij}(t)=P\{X_{s+t}=j\mid X_s=i\}.
\]

Тогда

\[
P(0)=I,
\qquad
p_{ij}(t)\ge0,
\qquad
\sum_jp_{ij}(t)=1,
\]

и

\[
P(t+s)=P(t)P(s).
\]

Это непрерывновременной аналог формулы \(P^{m+n}=P^mP^n\).

\subsection{5. Примеры}\label{ux43fux440ux438ux43cux435ux440ux44b}

\subsubsection{Пример 1. Двухсостоянийная
цепь}\label{ux43fux440ux438ux43cux435ux440-1.-ux434ux432ux443ux445ux441ux43eux441ux442ux43eux44fux43dux438ux439ux43dux430ux44f-ux446ux435ux43fux44c}

Пусть

\[
E=\{0,1\}.
\]

Матрица переходов:

\[
P=
\begin{pmatrix}
1-a & a\\
b & 1-b
\end{pmatrix},
\qquad
0\le a,b\le1.
\]

Из состояния \(0\) цепь переходит в \(1\) с вероятностью \(a\), а
остается в \(0\) с вероятностью \(1-a\). Из состояния \(1\) цепь
переходит в \(0\) с вероятностью \(b\), а остается в \(1\) с
вероятностью \(1-b\).

Строки суммируются в единицу:

\[
(1-a)+a=1,
\qquad
b+(1-b)=1.
\]

Если начальное распределение

\[
\pi_0=(\alpha,1-\alpha),
\]

то

\[
\pi_1=\pi_0P.
\]

\subsubsection{\texorpdfstring{Пример 2. Простое случайное блуждание на
\(\mathbb Z\)}{Пример 2. Простое случайное блуждание на \textbackslash mathbb Z}}\label{ux43fux440ux438ux43cux435ux440-2.-ux43fux440ux43eux441ux442ux43eux435-ux441ux43bux443ux447ux430ux439ux43dux43eux435-ux431ux43bux443ux436ux434ux430ux43dux438ux435-ux43dux430-mathbb-z}

Пусть

\[
E=\mathbb Z,
\]

и из состояния \(i\) цепь переходит в \(i+1\) с вероятностью \(p\), а в
\(i-1\) с вероятностью \(q=1-p\):

\[
p_{i,i+1}=p,
\qquad
p_{i,i-1}=q.
\]

Все остальные переходы равны нулю:

\[
p_{ij}=0,
\qquad
j\ne i+1,\ i-1.
\]

Это переходная матрица случайного блуждания из Билета 19. Она
бесконечная, но каждая строка имеет сумму

\[
p+q=1.
\]

\subsubsection{Пример 3. Поглощающее
состояние}\label{ux43fux440ux438ux43cux435ux440-3.-ux43fux43eux433ux43bux43eux449ux430ux44eux449ux435ux435-ux441ux43eux441ux442ux43eux44fux43dux438ux435}

Пусть

\[
E=\{0,1,2\}.
\]

Состояние \(0\) поглощающее:

\[
p_{00}=1,
\qquad
p_{01}=p_{02}=0.
\]

Например,

\[
P=
\begin{pmatrix}
1 & 0 & 0\\
1/3 & 1/3 & 1/3\\
1/2 & 0 & 1/2
\end{pmatrix}.
\]

Если цепь попала в \(0\), то она уже не выходит из \(0\). Это видно по
первой строке матрицы.

\subsubsection{Пример 4. Детерминированная
динамика}\label{ux43fux440ux438ux43cux435ux440-4.-ux434ux435ux442ux435ux440ux43cux438ux43dux438ux440ux43eux432ux430ux43dux43dux430ux44f-ux434ux438ux43dux430ux43cux438ux43aux430}

Пусть есть отображение

\[
T:E\to E.
\]

Если из состояния \(x\) система всегда переходит в \(T(x)\), то

\[
P(x,A)=\mathbf 1_A(Tx).
\]

В счетном случае:

\[
p_{ij}=
\begin{cases}
1, & j=T(i),\\
0, & j\ne T(i).
\end{cases}
\]

Это тоже марковское ядро. Случайность здесь может быть только в
начальном состоянии.

\subsubsection{Пример 5. Ядро с
плотностью}\label{ux43fux440ux438ux43cux435ux440-5.-ux44fux434ux440ux43e-ux441-ux43fux43bux43eux442ux43dux43eux441ux442ux44cux44e}

Пусть

\[
E=\mathbb R.
\]

Если при текущем состоянии \(x\) следующее состояние имеет плотность
\(p(x,y)\) по \(y\), то

\[
P(x,A)=\int_A p(x,y)\,dy.
\]

Для корректности нужно:

\[
p(x,y)\ge0,
\qquad
\int_{\mathbb R}p(x,y)\,dy=1
\]

для каждого \(x\), и измеримость по \(x\). Это непрерывный аналог
переходной матрицы.

\subsection{6. Схемы доказательств /
обоснований}\label{ux441ux445ux435ux43cux44b-ux434ux43eux43aux430ux437ux430ux442ux435ux43bux44cux441ux442ux432-ux43eux431ux43eux441ux43dux43eux432ux430ux43dux438ux439}

\textbf{Почему строки матрицы суммируются в единицу.}

Фиксируем текущее состояние \(i\). Тогда события

\[
\{X_{n+1}=j\},
\qquad
j\in E,
\]

образуют полную группу возможных значений следующего состояния. Поэтому

\[
\sum_jp_{ij}
=
\sum_jP\{X_{n+1}=j\mid X_n=i\}
=
1.
\]

\textbf{Почему \(\pi_{n+1}=\pi_nP\).}

Для каждого \(j\):

\[
\pi_{n+1}(j)
=
P\{X_{n+1}=j\}.
\]

Разбиваем по значениям \(X_n\):

\[
P\{X_{n+1}=j\}
=
\sum_iP\{X_n=i,\ X_{n+1}=j\}.
\]

Далее:

\[
P\{X_n=i,\ X_{n+1}=j\}
=
P\{X_n=i\}P\{X_{n+1}=j\mid X_n=i\}.
\]

Значит

\[
\pi_{n+1}(j)=\sum_i\pi_n(i)p_{ij}.
\]

Это и есть \(j\)-я координата строки \(\pi_nP\).

\textbf{Доказательство Чепмена-Колмогорова в счетном случае.}

Нужно найти вероятность перехода из \(i\) в \(j\) за \(m+n\) шагов:

\[
p_{ij}^{(m+n)}
=
P\{X_{m+n}=j\mid X_0=i\}.
\]

Разбиваем по промежуточному состоянию \(k\) в момент \(n\):

\[
p_{ij}^{(m+n)}
=
\sum_k
P\{X_n=k,\ X_{m+n}=j\mid X_0=i\}.
\]

По правилу умножения условных вероятностей:

\[
P\{X_n=k,\ X_{m+n}=j\mid X_0=i\}
=
P\{X_n=k\mid X_0=i\}
P\{X_{m+n}=j\mid X_n=k,\ X_0=i\}.
\]

По марковскому свойству и однородности будущее после момента \(n\)
зависит от прошлого только через \(X_n=k\):

\[
P\{X_{m+n}=j\mid X_n=k,\ X_0=i\}
=
P\{X_m=j\mid X_0=k\}
=
p_{kj}^{(m)}.
\]

Поэтому

\[
p_{ij}^{(m+n)}
=
\sum_k p_{ik}^{(n)}p_{kj}^{(m)}.
\]

\textbf{Доказательство Чепмена-Колмогорова для ядра.}

В общем пространстве состояний сумма по \(k\) заменяется интегралом по
промежуточному состоянию \(y\):

\[
P^{(m+n)}(x,A)
=
\int_E P^{(n)}(x,dy)P^{(m)}(y,A).
\]

Это та же формула полной вероятности, записанная для условных
распределений.

\textbf{Почему \(\pi_0\) и \(P\) задают конечномерные распределения.}

Для траектории

\[
i_0,i_1,\ldots,i_n
\]

используем цепное правило:

\[
P\{X_0=i_0,\ldots,X_n=i_n\}
=
P\{X_0=i_0\}
\prod_{k=1}^n
P\{X_k=i_k\mid X_0=i_0,\ldots,X_{k-1}=i_{k-1}\}.
\]

По марковскому свойству:

\[
P\{X_k=i_k\mid X_0=i_0,\ldots,X_{k-1}=i_{k-1}\}
=
p_{i_{k-1}i_k}.
\]

Значит

\[
P\{X_0=i_0,\ldots,X_n=i_n\}
=
\pi_0(i_0)p_{i_0i_1}\cdots p_{i_{n-1}i_n}.
\]

\subsection{7. Что написать на
доске}\label{ux447ux442ux43e-ux43dux430ux43fux438ux441ux430ux442ux44c-ux43dux430-ux434ux43eux441ux43aux435}

Переходное ядро:

\[
P(x,A)=P\{X_{n+1}\in A\mid X_n=x\}.
\]

Условия ядра:

\[
P(x,\cdot)\ \text{-- вероятность},
\qquad
P(\cdot,A)\ \text{-- измерима}.
\]

Матрица переходов:

\[
p_{ij}=P\{X_{n+1}=j\mid X_n=i\},
\qquad
p_{ij}\ge0,
\qquad
\sum_jp_{ij}=1.
\]

Эволюция распределения:

\[
\pi_{n+1}=\pi_nP,
\qquad
\pi_n=\pi_0P^n.
\]

Чепмен-Колмогоров:

\[
P^{(m+n)}(x,A)
=
\int_E P^{(n)}(x,dy)P^{(m)}(y,A).
\]

В счетном случае:

\[
p_{ij}^{(m+n)}
=
\sum_k p_{ik}^{(n)}p_{kj}^{(m)}.
\]

Формула траектории:

\[
P\{X_0=i_0,\ldots,X_n=i_n\}
=
\pi_0(i_0)p_{i_0i_1}\cdots p_{i_{n-1}i_n}.
\]

\subsection{8. Типичные дополнительные
вопросы}\label{ux442ux438ux43fux438ux447ux43dux44bux435-ux434ux43eux43fux43eux43bux43dux438ux442ux435ux43bux44cux43dux44bux435-ux432ux43eux43fux440ux43eux441ux44b}

\textbf{Что такое марковское ядро?}

Это функция \(P(x,A)\), которая при фиксированном \(x\) является
вероятностной мерой по \(A\), а при фиксированном \(A\) является
измеримой функцией по \(x\). Она задает условный закон следующего
состояния при текущем состоянии \(x\).

\textbf{Чем переходная матрица отличается от распределения цепи?}

Распределение \(\pi_n\) говорит, где цепь находится в момент \(n\).
Переходная матрица \(P\) говорит, как цепь переходит из одного состояния
в другое. Формула связи:

\[
\pi_{n+1}=\pi_nP.
\]

\textbf{Почему матрица переходов стохастическая именно по строкам?}

Потому что строка с номером \(i\) - это условное распределение
следующего состояния при условии \(X_n=i\):

\[
(p_{ij})_j.
\]

Сумма вероятностей всех возможных следующих состояний равна единице:

\[
\sum_jp_{ij}=1.
\]

\textbf{Что изменится при столбцовом соглашении?}

Если распределения записывать столбцами, то формула эволюции станет

\[
\pi_{n+1}=P^T\pi_n
\]

при той же строковой стохастичности. Либо можно использовать
столбцово-стохастические матрицы, но тогда нужно явно менять соглашение.
В этом конспекте используется строковое соглашение.

\textbf{Как получить переход за два шага?}

Нужно просуммировать по промежуточному состоянию:

\[
p_{ij}^{(2)}=\sum_kp_{ik}p_{kj}.
\]

Это элемент матрицы \(P^2\).

\textbf{Как формулируется уравнение Чепмена-Колмогорова?}

Для ядер:

\[
P^{(m+n)}(x,A)
=
\int_E P^{(n)}(x,dy)P^{(m)}(y,A).
\]

Для матриц:

\[
p_{ij}^{(m+n)}
=
\sum_kp_{ik}^{(n)}p_{kj}^{(m)}.
\]

\textbf{Почему Чепмен-Колмогоров следует из полной вероятности?}

Переход за \(m+n\) шагов можно разбить по состоянию \(k\) в
промежуточный момент \(n\). Суммируем вероятность попасть из \(i\) в
\(k\) за \(n\) шагов и затем из \(k\) в \(j\) за \(m\) шагов.

\textbf{Что значит однородность?}

Однородность означает, что переходный механизм не зависит от времени:

\[
P\{X_{n+1}\in A\mid X_n=x\}=P(x,A)
\]

для всех \(n\). В неоднородном случае нужно писать \(P_n(x,A)\).

\textbf{Можно ли задать марковскую цепь только матрицей переходов?}

Не полностью. Матрица задает переходный механизм, но еще нужно начальное
распределение \(\pi_0\). Вместе \(\pi_0\) и \(P\) задают все
конечномерные распределения.

\textbf{Почему случайное блуждание является примером переходной
матрицы?}

Потому что при известном текущем положении \(i\) следующее положение
зависит только от нового шага:

\[
p_{i,i+1}=p,
\qquad
p_{i,i-1}=q.
\]

Прошлая траектория не нужна для описания следующего перехода.

\subsection{9. Типичные
ошибки}\label{ux442ux438ux43fux438ux447ux43dux44bux435-ux43eux448ux438ux431ux43aux438}

\begin{itemize}
\tightlist
\item
  Делать стохастическими столбцы без оговорки, а потом использовать
  строковую формулу \(\pi_{n+1}=\pi_nP\).
\item
  Путать переходную матрицу \(P\) и распределение \(\pi_n\).
\item
  Забывать начальное распределение: одна матрица \(P\) не задает закон
  цепи полностью.
\item
  Писать \(p_{ij}^{(2)}=p_{ij}^2\). Правильно:
\end{itemize}

\[
p_{ij}^{(2)}=\sum_kp_{ik}p_{kj}.
\]

\begin{itemize}
\tightlist
\item
  В уравнении Чепмена-Колмогорова забывать сумму или интеграл по
  промежуточному состоянию.
\item
  Путать одношаговую вероятность \(p_{ij}\) и \(n\)-шаговую вероятность
  \(p_{ij}^{(n)}\).
\item
  Называть любую неотрицательную матрицу переходной. Нужна нормировка
  строк.
\item
  Путать однородность цепи с одинаковостью распределений \(X_n\).
  Однородность говорит о переходах, а не о том, что \(\pi_n\) постоянно.
\item
  Для общего пространства состояний забывать измеримость
  \(x\mapsto P(x,A)\).
\end{itemize}

\subsection{10. Связи с другими
билетами}\label{ux441ux432ux44fux437ux438-ux441-ux434ux440ux443ux433ux438ux43cux438-ux431ux438ux43bux435ux442ux430ux43cux438}

\begin{itemize}
\tightlist
\item
  Билет 05: переходное ядро требует измеримости по состоянию.
\item
  Билет 06: переходные функции задают условные распределения будущего
  состояния.
\item
  Билет 09: формулы с ядрами и распределениями опираются на
  интегрирование по произведениям и полную вероятность.
\item
  Билет 10: переходная функция является регулярным условным
  распределением.
\item
  Билет 15: по начальному распределению и переходным ядрам строятся
  согласованные конечномерные распределения.
\item
  Билет 16: марковская цепь - частный случай случайной
  последовательности.
\item
  Билет 19: случайное блуждание задается простой переходной матрицей.
\item
  Билет 23: переходные матрицы используются для определения и анализа
  марковских цепей.
\item
  Билет 24: для стационарных марковских цепей переходы влияют на
  корреляции и эргодичность.
\end{itemize}

\subsection{11. Минимум на
удовлетворительно}\label{ux43cux438ux43dux438ux43cux443ux43c-ux43dux430-ux443ux434ux43eux432ux43bux435ux442ux432ux43eux440ux438ux442ux435ux43bux44cux43dux43e}

Нужно сказать:

\begin{enumerate}
\def\labelenumi{\arabic{enumi}.}
\tightlist
\item
  Переходная функция:
\end{enumerate}

\[
P(x,A)=P\{X_{n+1}\in A\mid X_n=x\}.
\]

\begin{enumerate}
\def\labelenumi{\arabic{enumi}.}
\setcounter{enumi}{1}
\tightlist
\item
  В счетном случае:
\end{enumerate}

\[
p_{ij}=P\{X_{n+1}=j\mid X_n=i\}.
\]

\begin{enumerate}
\def\labelenumi{\arabic{enumi}.}
\setcounter{enumi}{2}
\tightlist
\item
  Матрица переходов стохастическая:
\end{enumerate}

\[
p_{ij}\ge0,
\qquad
\sum_jp_{ij}=1.
\]

\begin{enumerate}
\def\labelenumi{\arabic{enumi}.}
\setcounter{enumi}{3}
\tightlist
\item
  Распределения обновляются по формуле:
\end{enumerate}

\[
\pi_{n+1}=\pi_nP.
\]

\begin{enumerate}
\def\labelenumi{\arabic{enumi}.}
\setcounter{enumi}{4}
\tightlist
\item
  Чепмен-Колмогоров:
\end{enumerate}

\[
p_{ij}^{(m+n)}
=
\sum_kp_{ik}^{(n)}p_{kj}^{(m)}.
\]

\subsection{12. Минимум на
хорошо}\label{ux43cux438ux43dux438ux43cux443ux43c-ux43dux430-ux445ux43eux440ux43eux448ux43e}

Помимо минимума, нужно:

\begin{itemize}
\tightlist
\item
  объяснить, что \(P(x,\cdot)\) - мера, а \(P(\cdot,A)\) - измеримая
  функция;
\item
  различать однородные и неоднородные переходы;
\item
  вывести \(\pi_{n+1}=\pi_nP\) через полную вероятность;
\item
  вывести Чепмена-Колмогорова через промежуточное состояние;
\item
  привести пример: двухсостоянийная цепь, блуждание или поглощающее
  состояние.
\end{itemize}

\subsection{13. Что добавить для
отлично}\label{ux447ux442ux43e-ux434ux43eux431ux430ux432ux438ux442ux44c-ux434ux43bux44f-ux43eux442ux43bux438ux447ux43dux43e}

Для сильного ответа стоит добавить:

\begin{itemize}
\tightlist
\item
  формулу для общего ядра:
\end{itemize}

\[
\mu P(A)=\int_E P(x,A)\,\mu(dx);
\]

\begin{itemize}
\tightlist
\item
  формулу конечномерных распределений:
\end{itemize}

\[
P\{X_0=i_0,\ldots,X_n=i_n\}
=
\pi_0(i_0)p_{i_0i_1}\cdots p_{i_{n-1}i_n};
\]

\begin{itemize}
\tightlist
\item
  связь с теоремой Колмогорова: строки суммируются в единицу, поэтому
  конечномерные распределения согласованы;
\item
  замечание про непрерывное время:
\end{itemize}

\[
P(t+s)=P(t)P(s);
\]

\begin{itemize}
\tightlist
\item
  предупреждение про строковое и столбцовое соглашения.
\end{itemize}

\subsection{14. Устная версия ответа на 5
минут}\label{ux443ux441ux442ux43dux430ux44f-ux432ux435ux440ux441ux438ux44f-ux43eux442ux432ux435ux442ux430-ux43dux430-5-ux43cux438ux43dux443ux442}

Переходные функции и матрицы описывают вероятностный механизм движения
марковской системы. Пусть процесс \(X_n\) принимает значения в
пространстве состояний \((E,\mathcal E)\). Переходная функция задается
формулой

\[
P(x,A)=P\{X_{n+1}\in A\mid X_n=x\}.
\]

При фиксированном \(x\) это вероятностная мера по \(A\), а при
фиксированном \(A\) - измеримая функция по \(x\). Поэтому переходное
ядро можно применять к распределению \(\mu\):

\[
\mu P(A)=\int_E P(x,A)\,\mu(dx).
\]

Это распределение следующего состояния, если текущее состояние имеет
распределение \(\mu\).

В счетном пространстве состояний ядро задается матрицей

\[
p_{ij}=P\{X_{n+1}=j\mid X_n=i\}.
\]

Матрица стохастическая:

\[
p_{ij}\ge0,
\qquad
\sum_jp_{ij}=1.
\]

Здесь строки соответствуют текущему состоянию, а столбцы - следующему.
Если распределения записывать строками, то

\[
\pi_{n+1}=\pi_nP.
\]

Покомпонентно:

\[
\pi_{n+1}(j)=\sum_i\pi_n(i)p_{ij}.
\]

Переход за несколько шагов получается умножением матриц:

\[
P^{(n)}=P^n.
\]

Главная формула - Чепмен-Колмогоров:

\[
p_{ij}^{(m+n)}
=
\sum_kp_{ik}^{(n)}p_{kj}^{(m)}.
\]

Она получается из формулы полной вероятности: чтобы попасть из \(i\) в
\(j\) за \(m+n\) шагов, разбиваем траектории по промежуточному состоянию
\(k\) через \(n\) шагов.

Для общего пространства состояний сумма заменяется интегралом:

\[
P^{(m+n)}(x,A)
=
\int_E P^{(n)}(x,dy)P^{(m)}(y,A).
\]

Пример - случайное блуждание на \(\mathbb Z\):

\[
p_{i,i+1}=p,
\qquad
p_{i,i-1}=q,
\qquad
p+q=1.
\]

Все остальные переходы равны нулю. Это марковский переходный механизм:
чтобы узнать следующее положение, достаточно знать текущее.

В конце важно сказать, что переходная матрица не заменяет начальное
распределение. Закон цепи задается парой \(\pi_0\) и \(P\), а
конечномерные вероятности имеют вид

\[
P\{X_0=i_0,\ldots,X_n=i_n\}
=
\pi_0(i_0)p_{i_0i_1}\cdots p_{i_{n-1}i_n}.
\]

\subsection{15. Если осталось 2 минуты перед
экзаменом}\label{ux435ux441ux43bux438-ux43eux441ux442ux430ux43bux43eux441ux44c-2-ux43cux438ux43dux443ux442ux44b-ux43fux435ux440ux435ux434-ux44dux43aux437ux430ux43cux435ux43dux43eux43c}

Запомнить:

\[
P(x,A)=P\{X_{n+1}\in A\mid X_n=x\}.
\]

В счетном случае:

\[
p_{ij}=P\{X_{n+1}=j\mid X_n=i\},
\qquad
p_{ij}\ge0,
\qquad
\sum_jp_{ij}=1.
\]

Распределение:

\[
\pi_{n+1}=\pi_nP.
\]

\(n\)-шаговые переходы:

\[
P^{(n)}=P^n.
\]

Чепмен-Колмогоров:

\[
p_{ij}^{(m+n)}
=
\sum_kp_{ik}^{(n)}p_{kj}^{(m)}.
\]

Главная ловушка: строки, а не столбцы, являются условными
распределениями следующего состояния при фиксированном текущем
состоянии.

\newpage

\end{document}
